\section{Background and Related Work}
\label{sec:related}

\subsection{Substandard medicines and India's alerts}
WHO estimates that 10.5\% of medical products in low- and middle-income countries fail quality tests \cite{who2017oneinten}, and a meta-analysis put the pooled prevalence at 13.6\% \cite{ozawa2018prevalence}. A later review found that most failures are substandard rather than falsified \cite{ozawa2022characterizing}, which matters here because a substandard batch from a licensed maker looks genuine in every respect. Prevalence estimates depend strongly on study design \cite{mcmanus2020systematic}, and a 2025 scoping review found that evidence on harm is scarce and that interventions aimed at consumers are rare \cite{pal2025assessing}. In India, laboratories tested 116,323 samples in 2024--25 and found 3,104 not of standard quality and 245 spurious or adulterated \cite{roy2025fiveyears}. CDSCO's recall guidelines carry a recall to consumers only for Class I defects; Class II and III recalls stop at wholesalers and retailers \cite{cdsco2024recall}, and the 2012 version states that auditing and accountability for recalls were not in place \cite{cdsco2012recall}. A person therefore learns about a flagged batch only by reading the monthly alerts, which moved to an online viewer in August 2025 \cite{cdsco2025nsqlink}.

\subsection{Authentication, screening and existing tools}
India's Schedule H2 QR code is assigned per batch by each manufacturer and is mandatory for 300 brands \cite{cdsco2023faqqr}. WHO's traceability policy notes that verifying an identifier confirms the code, not the medicine inside the package, and that rogue verification services have answered every query positively \cite{who2021traceability}. Field screening devices range from US\$3 to over US\$60,000, and only 6 of 41 had been tested in the field \cite{vickers2018field,kovacs2014technologies}; in one field trial, handheld spectrometers silently chose the wrong reference library in 20 to 27\% of scans \cite{caillet2021comparative}. A review of 25 medicine-verification apps found almost no accuracy evidence and none that checked a regulator's quality list \cite{ciapponi2021mobile}. Image-based inspection of package design \cite{zakaria2025counterfeit,garciacotte2024deep} targets falsified packs, not batches that failed a laboratory test.

Several open tools built on CDSCO's data appeared in 2025 and 2026. Blacklisted offers typed batch lookup with an email digest \cite{blacklisted2026}. Nazar reads QR codes and the printed batch line from a photo, guards against look-alike characters by returning ``can't tell'' with a suggestion, and re-checks saved items when the user visits \cite{rastogi2026nazar}. SafeDose marks a batch hit with a different drug name as ``caution'' \cite{shaurya2026safedose}, and RecallOps India matches alerts against distributor stock using the batch, the product and one more field, with a human review queue \cite{ibnweb32026recallops}. Among the tools we found (searched October 2026), none asks the person to resolve a doubtful character during a check, corroborates product, brand, dates and maker on a photographed consumer pack, or runs fully offline. StripCheck adds these, together with an independent second reader and a verifier agent.

\subsection{Reading codes on packs}
Batch and expiry codes are hard to read. On 27 pharmaceutical pack images, a region-based detector reached an F1 score of 80.9 against 47.1 for Tesseract \cite{koponen2024novel}, while expiry dates on food packages are read with about 97\% accuracy \cite{seker2022generalized}. StripCheck uses the PP-OCR family \cite{du2020ppocr,du2021ppocrv2,li2022ppocrv3}. Confidence scores of modern networks are poorly calibrated \cite{guo2017calibration}, and vision-language models are no better at signalling when to abstain on text \cite{yao2025reading}; on degraded text they fall back on language priors \cite{he2025seeing}, so a model's fluent reading of a code is not evidence that the code is right. Contrast enhancement with CLAHE did not improve licence-plate reading \cite{tavares2024comparison}, consistent with our observation that an enhanced re-read repeats a confident misread.

\subsection{Post-OCR correction and hotlist matching}
Noisy-channel models rank candidate corrections by how likely the observed characters are given each candidate \cite{kernighan1990spelling}. A learned error model raised one-best accuracy from 89.5\% to 95.0\% \cite{brill2000improved}, and a channel model over a closed vocabulary cut OCR word error by 70\% \cite{kolak2002ocr,kolak2003generative}. Correction remains harder than detection: in the ICDAR 2019 benchmark, detection reached 95\% but the best correction improved text by only 44\% \cite{rigaud2019icdar,nguyen2021survey}. Routing only the words on which two OCR engines disagree to people raised accuracy to 99.1\% \cite{vonahn2008recaptcha}, the pattern StripCheck follows with its second reader and character questions. Licence-plate recognition is the closest operational analogue. Its confusable pairs (B/8, D/0, I/1, S/5, Z/2) match those on batch codes \cite{laroca2021efficient,du2013alpr}; in Seattle about one hotlist hit in five was a misread, and each hit required an explicit accept-or-misread decision \cite{seattleoig2023alpr}; and a misread of 3 as 7 led to a stop at gunpoint, in which the court stressed visual verification before acting \cite{green2014ninth}. StripCheck's confirmation gate applies the same rule to medicine alarms.

\subsection{Record linkage}
Fellegi and Sunter sort record pairs into links, non-links and possible links, which ``will require expensive manual linkage operations'' \cite{fellegi1969theory}. StripCheck maps these zones to an alarm, \textsc{not\_listed} and a question or abstention. Errors in a blocking key lose true matches, and exact blocking on keys with typos can miss about a quarter of them, so practical systems use several blocking passes \cite{winkler2006overview,christen2012survey}; confusion variants of the batch key play that role. Candidate generation can be reported separately from final accuracy \cite{papadakis2020survey}. Language models match entities nearly as well as fine-tuned matchers but are sensitive to the prompt \cite{peeters2025entity}, and small fine-tuned models perform on par with large prompted ones \cite{zhang2025deep}, which supports a deterministic matcher with a small model only for judging product names. Indian brand names collide across unrelated molecules \cite{neelakantan2024lookalike}, so a brand match alone never raises an alarm.

\subsection{LLM agents and external verification}
Agents that interleave reasoning and tool use \cite{yao2023react,schick2023toolformer} and that critique or refine their own output \cite{shinn2023reflexion,madaan2023selfrefine,gou2024critic} are now common, but models cannot reliably correct their own reasoning without external feedback \cite{huang2024cannot}, and planning works best when an external verifier checks the model's proposals \cite{kambhampati2024llmmodulo}. Reliability is low even for strong models: agents completed under half of the tasks in $\tau$-bench \cite{yao2024taubench}, and open models trail commercial ones \cite{liu2024agentbench}. Small models suit narrow agent tasks \cite{belcak2025slm}, though strict output formats can reduce reasoning quality \cite{tam2024speak}, and agent comparisons should report cost and use simple baselines \cite{kapoor2024agents}. StripCheck therefore lets a small model choose only among allowed actions at branch points, keeps a deterministic planner as fallback, and checks every draft with hard rules, using LangGraph interrupts for questions to the user \cite{langchain2026interrupts}.

\subsection{Abstention, oversight and risk communication}
The reject option \cite{chow1970reject} and selective classification \cite{geifman2017selective} formalise answering only when confident. Language models rarely abstain on medical questions \cite{cocchieri2026abstain} and often lack support for their citations \cite{gao2023alce}, so StripCheck's abstention and cited verdicts come from rules and templates. People over-rely on automated advice regardless of training \cite{parasuraman2010complacency} and accept explained advice whether or not it is correct; cognitive forcing functions reduce this over-reliance \cite{bansal2021whole,bucinca2021trust}, which is why StripCheck's confirmation asks the person to read specific characters from the pack. Adding a residual-risk sentence to a negative test result raised correct understanding from 54\% to 89\% \cite{batteux2022residual}, which supports the wording of \textsc{not\_listed}. Symptom checkers show how consumer tools err \cite{semigran2015symptom}, and WHO's guidance on large multimodal models applies to the language-model parts \cite{who2024lmm}. For Indian users, self-medication is common \cite{rashid2020selfmed}, pharmacies sell many medicines without prescription \cite{chandran2022antibiotics}, and text-only interfaces fail low-literacy users while graphics and voice help \cite{medhi2007textfree,medhi2011mobile}.

\subsection{Recalls and data freshness}
Consumer-facing recall systems elsewhere publish how current they are: US vehicle recall lookups must be updated at least every seven days and show the last update \cite{ecfr2026vin}, a convention StripCheck adopts for its alert index. Recalls often arrive late, with 34\% of 50 substandard-drug recalls in Nepal issued after 15 months on the market \cite{barma2025recall}, and recall letters alone changed little in one clinical setting \cite{morizio2021veterans}. Matching recalls to a person's own medicines is feasible when the dispensed lot is known \cite{gadgil2026automating}, which is what StripCheck's medicine cabinet records.
